%==============================================================
%  STYLE B — Minimaliste moderne (Noir + Rouge cerise)
%==============================================================
\documentclass[a4paper,11pt]{article}
\usepackage[T1]{fontenc}
\usepackage[french]{babel}
\usepackage{lmodern}
\usepackage{microtype}
\usepackage{setspace}
\setstretch{1.25}
\usepackage{amsmath, amssymb}
\usepackage[dvipsnames]{xcolor}
\usepackage{tikz}
\usepackage[most]{tcolorbox}
\usepackage[framemethod=tikz]{mdframed}
\usepackage[top=2.5cm, bottom=2.5cm, left=3cm, right=2.5cm]{geometry}
\usepackage{enumitem}
\usepackage{titlesec}

%=== PALETTE — Noir & Rouge cerise ==========================
\definecolor{nearblack}{HTML}{1A1A1A}
\definecolor{crimson}{HTML}{B5121B}
\definecolor{lightgray}{HTML}{F4F4F4}
\definecolor{midgray}{HTML}{888888}
\definecolor{accentvt}{HTML}{1A6B3C}

%=== BOÎTE THÉORÈME ==========================================
\tcbset{
  thmbox/.style={
    enhanced,
    colback=white,
    colframe=nearblack,
    boxrule=0.8pt,
    arc=0pt,
    left=12pt, right=10pt, top=8pt, bottom=9pt,
    before skip=12pt, after skip=14pt,
    borderline west={3pt}{0pt}{crimson},
  }
}

%=== BARRES LATÉRALES ========================================
\mdfdefinestyle{proofbar}{
  topline=false, bottomline=false, rightline=false,
  linecolor=midgray, linewidth=2pt,
  innerleftmargin=14pt, innerrightmargin=0pt,
  innertopmargin=3pt, innerbottommargin=7pt,
  backgroundcolor=white,
  skipabove=4pt, skipbelow=10pt,
}
\mdfdefinestyle{exbar}{
  topline=false, bottomline=false, rightline=false,
  linecolor=accentvt, linewidth=2pt,
  innerleftmargin=14pt, innerrightmargin=0pt,
  innertopmargin=3pt, innerbottommargin=7pt,
  backgroundcolor=lightgray,
  skipabove=4pt, skipbelow=8pt,
}
\mdfdefinestyle{exobar}{
  topline=false, bottomline=false, rightline=false,
  linecolor=crimson, linewidth=2pt,
  innerleftmargin=14pt, innerrightmargin=0pt,
  innertopmargin=3pt, innerbottommargin=8pt,
  backgroundcolor=white,
  skipabove=4pt, skipbelow=6pt,
}

%=== TITRE DE LEÇON ==========================================
\newcommand{\lecontitle}[2]{%
  \thispagestyle{empty}
  \vspace*{6pt}
  \noindent
  \begin{tikzpicture}
    \fill[crimson] (0,0) rectangle (3pt, 1.6cm);
    \node[anchor=west] at (10pt, 0.8cm)
      {\fontsize{22}{28}\selectfont\bfseries\color{nearblack} #1};
  \end{tikzpicture}
  \par\vspace{2pt}
  {\large\color{midgray} #2}
  \par\vspace{20pt}
}

%=== NUMÉRO DE SECTION =======================================
\newcommand{\secnum}[3]{%
  \vspace{18pt}
  \noindent
  {\sffamily\large\bfseries\color{midgray} 0#2\enspace}%
  {\sffamily\large\bfseries\color{nearblack} #3}%
  \par\vspace{2pt}
  \noindent\textcolor{nearblack!15}{\rule{\textwidth}{0.5pt}}
  \vspace{6pt}
}

%=== RUBRIQ ==================================================
\newcommand{\rubriq}[1]{%
  \vspace{8pt}
  \noindent{\bfseries\color{crimson} #1.}\enspace
}

%=== SÉPARATEUR ==============================================
\newcommand{\decsep}{%
  \vspace{12pt}
  \begin{center}
    \textcolor{nearblack!20}{\rule{6cm}{0.5pt}}
  \end{center}
  \vspace{6pt}
}

%=== PIED DE LEÇON ===========================================
\newcommand{\leconfooter}[1]{%
  \vfill
  \noindent\textcolor{nearblack!15}{\rule{\textwidth}{0.5pt}}
  \begin{center}
    \small\color{midgray}
    Fiche de cours \quad|\quad #1
  \end{center}
}

\setlist[enumerate]{label=\textcolor{crimson}{\arabic*.}, itemsep=4pt, topsep=6pt}
\setlist[itemize]{itemsep=4pt, topsep=6pt,
  label=\textcolor{crimson}{\textbullet}}

\pagestyle{empty}

\begin{document}

%--------------------------------------------------------------
\lecontitle{Théorème des valeurs intermédiaires}{Analyse réelle — Fonctions continues}

\secnum{}{1}{Énoncé}

\begin{tcolorbox}[thmbox]
  \textbf{Théorème (TVI).}\enspace
  Soit $f : [a,b] \to \mathbb{R}$ continue. Pour tout $y$ compris entre
  $f(a)$ et $f(b)$, il existe $c \in [a,b]$ tel que $f(c) = y$.

  \medskip En particulier, si $f(a) < 0 < f(b)$, alors $f$ admet une
  racine dans $]a,b[$.
\end{tcolorbox}

\rubriq{Ingrédient clé}
La preuve repose sur la \textbf{complétude de $\mathbb{R}$} (propriété de
la borne supérieure) et la continuité de $f$.

\decsep

\secnum{}{2}{Preuve}

\begin{mdframed}[style=proofbar]
  \textit{Étape 1.} Posons $E = \{x \in [a,b] \mid f(x) \leq y\}$.
  L'ensemble $E$ est non vide (car $a \in E$) et majoré par $b$.

  \medskip
  \textit{Étape 2.} Posons $c = \sup E$. Par continuité de $f$ en $c$,
  $f(c) = y$.\hfill$\blacksquare$
\end{mdframed}

\decsep

\secnum{}{3}{Exemples et exercices}

\begin{mdframed}[style=exbar]
  \textbf{Exemple.}\enspace $f(x) = x^3 - 2x - 5$ vérifie $f(2) = -1 < 0$
  et $f(3) = 16 > 0$, donc $f$ admet une racine dans $]2,3[$.
\end{mdframed}

\vspace{6pt}

\begin{mdframed}[style=exobar]
  \begin{enumerate}
    \item Montrer que $e^x = 2-x$ a une solution réelle.
    \item Soit $f$ continue sur $[0,1]$ avec $f(0) = f(1) = 0$.
          Montrer que $f$ n'est pas nécessairement dérivable.
  \end{enumerate}
\end{mdframed}

\leconfooter{B.\ Bolzano (1817) — A.-L.\ Cauchy (1821)}

\end{document}
